I can prove Parts A–C exactly as stated. For Part D, however, the displayed formula shows that the conclusion \(J(x)<0\) is false as written; the correct conclusion is \(J(x)=0\) whenever the integral is well-defined. I will prove this rigorously.

The only external fact I use is Euler’s formula
\[
\cos u+i\sin u=e^{iu},
\]
from which, by replacing \(u\) with \(-u\) and \(u+\pi\), one gets
\[
\sin(-u)=-\sin u,\qquad \cos(-u)=\cos u,\qquad \sin(u+\pi)=-\sin u,\qquad \cos(u+\pi)=-\cos u.
\]
These are standard and are tabulated in DLMF. ([dlmf.nist.gov](https://dlmf.nist.gov/4.14.E9))

Fix \(x\in(0,1)\), and write
\[
r:=\sqrt{1-x^2},\qquad y:=r\sin\psi,\qquad z:=r\cos\psi.
\]
Then
\[
x^2+y^2+z^2=1,
\]
and the given formula is
\[
h(x,\psi)= -\frac{x\,y\,(3+2z^2-x^2)}{9\pi\,P_A(x,z)},
\]
where
\[
P_A(x,z)=\bigl((1-az)^2-b^2x^2\bigr)\bigl((1+az)^2-b^2x^2\bigr),
\qquad a=\sqrt{\tfrac23},\quad b=\tfrac1{\sqrt3}.
\]

## Part A: \(h(x,-\psi)=-h(x,\psi)\)

Since
\[
y(x,-\psi)=r\sin(-\psi)=-r\sin\psi=-y(x,\psi),
\]
and
\[
z(x,-\psi)=r\cos(-\psi)=r\cos\psi=z(x,\psi),
\]
we have
\[
\begin{aligned}
h(x,-\psi)
&=-\frac{x\,y(x,-\psi)\,\bigl(3+2z(x,-\psi)^2-x^2\bigr)}
{9\pi\,P_A\bigl(x,z(x,-\psi)\bigr)}\\
&=-\frac{x(-y)\,(3+2z^2-x^2)}{9\pi P_A(x,z)}\\
&=\frac{x\,y\,(3+2z^2-x^2)}{9\pi P_A(x,z)}\\
&=-h(x,\psi).
\end{aligned}
\]
Thus \(h(x,\cdot)\) is odd in \(\psi\).

## Part B: \(h(x,\psi)<0\) for \(\psi\in(0,\pi/2)\)

In fact, we prove the stronger statement
\[
h(x,\psi)<0\qquad\text{for all }\psi\in(0,\pi).
\]

First, write
\[
Q_-(x,\psi):=(1-az)^2-b^2x^2,\qquad Q_+(x,\psi):=(1+az)^2-b^2x^2.
\]
Using \(a^2=\tfrac23\), \(b^2=\tfrac13\), and \(x^2+y^2+z^2=1\), we get
\[
\begin{aligned}
Q_-(x,\psi)
&=1-2az+\frac23 z^2-\frac13 x^2\\
&=1-2az+\frac23 z^2-\frac13(1-y^2-z^2)\\
&=(z-a)^2+\frac13 y^2,
\end{aligned}
\]
and similarly
\[
Q_+(x,\psi)=(z+a)^2+\frac13 y^2.
\]
Hence
\[
P_A(x,z)=Q_-(x,\psi)Q_+(x,\psi)>0
\]
whenever \(y\neq 0\). In particular, for \(\psi\in(0,\pi)\), we have \(y=r\sin\psi>0\), so \(P_A(x,z)>0\).

Next,
\[
3+2z^2-x^2
=2+(1-x^2)+2z^2
=2+y^2+3z^2>0.
\]

Therefore, for \(\psi\in(0,\pi)\),
\[
x>0,\qquad y>0,\qquad 3+2z^2-x^2>0,\qquad P_A(x,z)>0.
\]
So
\[
h(x,\psi)= -\frac{\text{positive}}{\text{positive}}<0.
\]
In particular,
\[
h(x,\psi)<0\qquad\text{for all }\psi\in(0,\pi/2).
\]

## Part C: \(h(x,\psi)\,h(x,\psi-\pi/2)<0\) for \(\psi\in(0,\pi/2)\)

Let \(\psi\in(0,\pi/2)\). Then
\[
\frac{\pi}{2}-\psi\in(0,\pi/2)
\]
and
\[
\psi-\frac{\pi}{2}=-\Bigl(\frac{\pi}{2}-\psi\Bigr).
\]
By Part A,
\[
h\Bigl(x,\psi-\frac{\pi}{2}\Bigr)
=-h\Bigl(x,\frac{\pi}{2}-\psi\Bigr).
\]
By Part B,
\[
h(x,\psi)<0
\qquad\text{and}\qquad
h\Bigl(x,\frac{\pi}{2}-\psi\Bigr)<0.
\]
Hence
\[
h\Bigl(x,\psi-\frac{\pi}{2}\Bigr)>0.
\]
Therefore
\[
h(x,\psi)\,h\Bigl(x,\psi-\frac{\pi}{2}\Bigr)<0
\qquad\text{for all }\psi\in(0,\pi/2).
\]

## Part D: the stated conclusion \(J(x)<0\) is not correct

The issue is that the function
\[
H_x(\psi):=h(x,\psi)\,h\Bigl(x,\psi-\frac{\pi}{2}\Bigr)
\]
is **not** \(\pi/2\)-periodic. It is \(\pi/2\)-**anti**-periodic.

Indeed, from
\[
\sin(\psi+\pi)=-\sin\psi,\qquad \cos(\psi+\pi)=-\cos\psi,
\]
we get
\[
y(\psi+\pi)=-y(\psi),\qquad z(\psi+\pi)=-z(\psi).
\]
Also,
\[
P_A(x,-z)=P_A(x,z),
\]
because the two factors in \(P_A\) are interchanged. Therefore
\[
\begin{aligned}
h(x,\psi+\pi)
&=-\frac{x(-y)(3+2z^2-x^2)}{9\pi P_A(x,-z)}\\
&=\frac{x\,y\,(3+2z^2-x^2)}{9\pi P_A(x,z)}\\
&=-h(x,\psi).
\end{aligned}
\]
So \(h\) is \(\pi\)-anti-periodic.

Now
\[
\begin{aligned}
H_x\Bigl(\psi+\frac{\pi}{2}\Bigr)
&=h\Bigl(x,\psi+\frac{\pi}{2}\Bigr)h(x,\psi)\\
&=h\Bigl(x,\Bigl(\psi-\frac{\pi}{2}\Bigr)+\pi\Bigr)h(x,\psi)\\
&=-h\Bigl(x,\psi-\frac{\pi}{2}\Bigr)h(x,\psi)\\
&=-H_x(\psi).
\end{aligned}
\]
Thus \(H_x\) is \(\pi/2\)-anti-periodic, not \(\pi/2\)-periodic.

Since \(w(x,\psi+\pi/2)=w(x,\psi)\), the full integrand
\[
I_x(\psi):=H_x(\psi)\,w(x,\psi)
\]
also satisfies
\[
I_x\Bigl(\psi+\frac{\pi}{2}\Bigr)=-I_x(\psi).
\]
Hence, whenever the integral is well-defined,
\[
\int_0^{2\pi} I_x(\psi)\,d\psi
=\int_0^{\pi/2}\sum_{k=0}^3 I_x\Bigl(\psi+k\frac{\pi}{2}\Bigr)\,d\psi
=0.
\]
So the correct conclusion is
\[
J(x)=0,
\]
not \(J(x)<0\).

Finally, from the factorization in Part B, \(P_A(x,z)=0\) can occur only if
\[
y=0\quad\text{and}\quad z=\pm a,
\]
which forces
\[
x^2=1-a^2=\frac13,\qquad x=\frac1{\sqrt3}=b.
\]
Thus for
\[
x\neq \frac1{\sqrt3},
\]
the integrand is continuous on \([0,2\pi]\), and the above cancellation is legitimate. At
\[
x=\frac1{\sqrt3},
\]
the denominator vanishes at \(\psi=0,\pi,2\pi\), so \(J(x)\) is not an ordinary integral without further interpretation.

## Conclusion

The rigorous conclusions from the explicit formula are:

\[
h(x,-\psi)=-h(x,\psi),
\]
\[
h(x,\psi)<0\quad\text{for }\psi\in(0,\pi/2)\ \text{(indeed for all }\psi\in(0,\pi)\text{)},
\]
\[
h(x,\psi)\,h\Bigl(x,\psi-\frac{\pi}{2}\Bigr)<0\quad\text{for }\psi\in(0,\pi/2).
\]

But Part D, as written, is false. The correct replacement is:
\[
J(x)=0
\]
whenever the integral is well-defined.